Część teoretyczna (3 rozdziały 10-20 stron max)
1.	Nowoczesne systemy aplikacyjne
a.	założenia Cloud Native Computing 
b.	Krajobraz Cloud Native Computing Foundation 
c.	architektura Twelve-Factor App 
2.	DevOps , Platform Engineering   i SRE  – krótkie omówienie nowoczesnych strategii pracy rozwoju i utrzymywania systemów
3.	Etapy wytwarzania oprogramowania wg. założeń DevOps – pętla ciągłego rozwoju 
4.	Jak DevOps i Cloud Native się uzupełniają?
a.	Podsumowanie i najważniejsze wnioski – zarys części praktycznej

Cześć praktyczna (szkic)
Implementacja aplikacji wykorzystującej narzędzia i koncepcje DevOps + Cloud Native
1.	Plan (opis na narzędzi do planowania zadań, przykłady rozplanowania funkcji systemu, integrtacja z systemem kontroli wersji)
a.	Jira
2.	Code (narzędzia użyte do wytwarzania oprogramowania, alternatywy, architektura aplikacji z uzględnieniem 12 factor app)
a.	System kontroli werjsji: GitHub
b.	IDE: VS Code
c.	Technologia backend: Java, Quarkus
d.	Technologia frontend: React
3.	Build (budowanie i wersjonowanie paczek)
a.	GitHub Actions – proces CI/CD zapisany w kodzie
b.	Gradle – budowanie backendu
c.	Npm/Webpack – budowanie frontendu
4.	Test (etap procesu CI/CD wbudowany w GitHub Actions)
a.	GitHub Actions
b.	JUnit
5.	Release (wydawanie paczek releasowych, zapisywanie kolejnych iteracji projektów) 
a.	Docker + DockerHub – publiczne repozytorium dla kontenerów
b.	Helm + ArtifactHub – publiczne repozytorium dla schematów Helm
6.	Deploy (infrastuktura wspierająca, klaster do uruchomienia kontenerów)
a.	Terraform – tworzenie zasobów na AWS wg. koncepcji Infrastructure-as-a-Code
b.	Ansible – konfiguracja serwerów z użyciem Configuration Management
c.	Kubernetes – środowisko uruchomieniowe aplikacji
7.	Operate (środowisko chmurowe, serwery, usługi wspierające)
a.	AWS Cloud Service
i.	RDS – bazy danych
ii.	S3 – zapis plików
iii.	Opensearch – zapis dokumentów
iv.	(inne usługi – lista może się zmienić w trakcie implementacji)
b.	Kubernetes
8.	Monitor + Feedback (monitorowanie systemu, aplikacji, dziennki logów, ślady aplikacji)
a.	Grafana – wizualizacja metryk
b.	Prometheus – zbieranie metryk
c.	Elasticsearch/Opensearch – zbieranie logów
d.	Jaeger lub Zipkin – zbieranie trace’ów
e.	(OpenTelemetry) – standard komunikacyjny do powyższych technologii
f.	OpsGenie – alerting 